\cotefile{19}
\scanpage{81}
adjoint à droite $h$, on aura la suite \unsure{étendue}

\diagram{\begin{tikzcd}[ampersand replacement=\&,column sep=small,row sep=small]
\deleted{\unsure{\varphi_{!}}} \& \deleted{\unsure{\varphi^*|\varphi}} \arrow[d,equal] \& \deleted{\unsure{\varphi_*}} \arrow[d,equal] \& \& \& \\
\& f_! \& f^* \arrow[d,equal] \arrow[rr,no head,to path={-- ([yshift=2mm]\tikztostart.north) -| (\tikztotarget.north)}] \& f_* \arrow[d,equal] \& f^! \arrow[d,equal] \& f_? \arrow[d,equal]\\
\& \& g_! \& g^* \arrow[d,equal] \& g_* \arrow[d,equal] \& g^! \arrow[d,equal]\\
\& \& \& h_! \& h^* \& h_*
\end{tikzcd}}

etc. On trouve ainsi le résultat (de proche en proche)

\uline{Proposition} Soit $\mathcal U$ un univers, \deleted{\unsure{$\mathbb N$}} contenant
un ensemble infini. Pour \unsure{tt|tout} entier $n$,
il existe deux $\mathcal U$-topos $\unsure{E|\mathcal E}$, $F$
(de la forme $\hat C$, $\hat C'$ si on veut, avec
$C$, $C'$ des petites catégories) et une
suite de $n$-foncteurs \inserted{\deleted{\illegible{1}}, dont aucun n'est pl. fid.,} $f_1,f_2,f_3\ \ldots,f_n$
alternativement de $E$ dans $F$ et de $F$ dans $E$,
chacun étant l'adjoint à droite de celui
qui précède. \unsure{Celà|Ceci} si on veut, tous les $f_{\unsure{i|j}}$
y inclus $=$ $f_1$ et $f_n$, étant des morphismes
de topos -- et \unsure{$=$} des morphismes \unsure{stricts|direct}].

\uline{Question}. Peut-on trouver une suite infinie
dans les deux directions ? Il suffirait de
trouver $f$ et $g$, tels que $g$ soit à la fois adjoint
à gauche et \unsure{,|à} droite de $f$, i.e. tel que les

% archive sheet number: 80
